
In this section, we provide background on reasoning, describe the key vulnerability introduced by reasoning being widely compatible and then how this can be exploited for scalable extraction of traces.

\subsection{Threat Model}
We assume a standard, unprivileged API adversary. The attacker does not require insider access to the provider’s infrastructure, cannot observe server-side states, and has no access to the proprietary model weights. Instead, the adversary operates entirely within the boundaries of standard API usage. We consider two attacker profiles:
\begin{itemize}
    \item \textbf{First-Party Attacker (Distillation \& Jailbreaking):} The adversary generates their own encrypted reasoning traces by querying a capable, safeguarded target model. They then exploit cross-model compatibility to replay these traces into a weaker, cheaper decoder model to deliberately bypass alignment guardrails or extract proprietary chains of thought. 
    \item \textbf{Third-Party Attacker (Secret Extraction \& Prompt Injection):} The adversary intercepts, scrapes, or otherwise acquires legitimate encrypted reasoning blobs generated by other users -- such as developers publishing raw agent session logs online. The attacker then leverages cross-user compatibility to act as a decryption oracle, uncovering sensitive data hidden in the victim's traces, or injecting maliciously crafted opaque blocks back into the victim's workflow.  
\end{itemize}
    
The common attack prerequisite is access to a compatible ``decoder'' model within the provider ecosystem.

\subsection{Encrypted Reasoning}
\label{sec:vuln_mechanism}

Chain-of-thought reasoning \citep{wei2022chain, jaech2024openai} allows modern LLMs to dynamically adjust test-time compute to reason over complex tasks prior to generating a user-directed response. These intermediate steps capture a model's problem-solving process, resulting in traces that are significantly more detailed and information-dense than the output provided to the user. As this internal monologue can expose the underlying reasoning methods as well as model training techniques, it is considered highly valuable for the developers of competing systems.

To protect from unauthorized extraction, most major API providers have deprecated the transmission of plaintext reasoning. Instead, modern APIs implement \emph{extended-thinking blocks}. When a model generates reasoning, the API returns a block where the human-readable component is either entirely hidden or heavily summarized. To avoid server-side storage, however, the actual chain-of-thought payload is still packaged into an opaque, \texttt{base64}-encoded signature or encrypted payload. This string functions as an Authenticated Encryption with Associated Data (AEAD) envelope, containing a header, which, depending on providers, can specify the model name, block type, version, and key ID along with a nonce, an authentication tag, and the ciphertext \citep{green2026encryptedreasoning}. The signature acts as associated data that is hashed into the Message Authentication Code (MAC), allowing the provider to verify and replay the block without storing it server-side. Depending on the API integration, this envelope is passed back to the API in consecutive calls using fields such as \texttt{signature} or \texttt{thinkingSignature}.

This design has three critical operational functions:
\begin{itemize}
    \item \textbf{Confidentiality.} \; By rendering the reasoning steps mathematically opaque, model providers block competitors from mass-harvesting explicit chain-of-thought data for model distillation and limit exposure of a potentially sensitive or harmful reasoning.
    \item \looseness=-1 \textbf{Integrity.} \; The AEAD envelope and MAC ensure that the intermediate reasoning cannot be maliciously altered by the user; any tampering invalidates the signature and can be rejected by the API, preventing manipulated reasoning from being used to steer the model.
    \item \textbf{Statelessness.} \; Rather than incurring overheads of storing reasoning states, providers leverage client-side storage. The client is required to hold the encrypted trace and pass it back in subsequent API calls to maintain the continuity within multi-turn session. 
\end{itemize}
As of July 2026 no LLM provider provides detailed description of the cryptographic mechanism used. Given the stateless design on the protocols, we can infer that reliance on client-side storage requires that the encrypted blobs are portable across contexts, users, and models. This enables functionality such as seamless model switching and automatic re-routing without discarding reasoning tokens. While model providers could limit this form of compatibility, and e.g. allow decryption only within the same conversation with the respective model, most APIs opt for a simpler strategy of enabling broader compatibility.


\subsection{Reasoning Compatibility}
\label{sec:compatibility}

To enable portability of reasoning traces across model calls, our experiments show that providers appear to be using a single global key to encrypt and authenticate every reasoning block. \citet{green2026encryptedreasoning} recognized this and showed that this allows clients to replay reasoning blocks out of order or across sessions.

For the purposes of our vulnerability analysis, we distinguish three forms of \textit{reasoning compatibility} of increasing permissiveness, each enabling a broader class of attacks.


\textbf{In- and cross-session compatibility.} \; A user can replay reasoning blocks in a different order than they were produced by the model. The user can also reuse blocks from earlier sessions in new requests. While this simplifies benign history editing and context truncation, it also allows attackers to fabricate conversation history (e.g., inserting compliant thoughts into an otherwise malicious exchange \citep{khalil2026hiddenthought}). In this work, we use this form of manipulation to perform reasoning extraction as discussed in \Cref{sec:extraction}.

\textbf{Cross-user compatibility.} \; With cross-user compatibility, a user can replay encrypted reasoning blocks, taken from another user's sessions. Combined with the above, this enables extraction of sensitive information such as API secrets by third parties (\Cref{sec:av_secret_extraction}), targeting, for instance, private information that has leaked into a model's reasoning \citep{green2025leakythoughtslargereasoning}.


\textbf{Cross-model compatibility.} \; With cross-model compatibility, a user can replay reasoning blocks produced by one model in a request to another. While this compatibility enables seamless model downgrades (e.g., from Opus to Sonnet or from Opus~4.8 to Opus~4.6), it also plays a central role in our attack method described in \Cref{sec:extraction}. In particular, it enables scalable distillation of a highly capable model reasoning without ever querying that model for the said reasoning directly (see \Cref{fig:teaser}). \Cref{tab:crossmodel} provides an overview of cross-model compatibility across key providers.


\begin{table}[h]
\centering
\caption{\textbf{Cross-model compatibility of encrypted reasoning.} As per July 2026. Row: the source model that produced the encrypted reasoning block; column: the target model receiving the injected reasoning. A \cmark \, indicates that, for this combination, the target model interacts with the injected thought.
\textbf{Claude:} the thinking traces of any model can be replayed by any other, except Fable 5's thoughts.
\textbf{GPT:} the GPT-5.6 series can replay the traces
of all earlier model generations. \textbf{Gemini:} the thinking traces of any model can
be replayed into any other.}
\label{tab:crossmodel}
\tiny
\setlength{\tabcolsep}{1.2pt}
\renewcommand{\arraystretch}{0.95}
\resizebox{\textwidth}{!}{%
\begin{tabular}{@{}l@{\;}cccccc@{\hspace{1.3em}}l@{\;}cccccc@{\hspace{1.3em}}l@{\;}cccccc@{}}
\toprule
\multicolumn{7}{c}{\scriptsize Claude} &
\multicolumn{7}{c}{\scriptsize GPT} &
\multicolumn{7}{c}{\scriptsize Gemini} \\
\cmidrule(r{0.9em}){1-7}\cmidrule(lr{0.9em}){8-14}\cmidrule(l){15-21}
\textbf{Source / Target} & \textbf{F5} & \textbf{O4.8} & \textbf{S5} & \textbf{S4.6} & \textbf{S4.5} & \textbf{H4.5} &
\textbf{Source / Target} & \textbf{5.6s} & \textbf{5.6t} & \textbf{5.6l} & \textbf{5} & \textbf{5-m} & \textbf{5-n} &
\textbf{Source / Target} & \textbf{3.1P} & \textbf{3P} & \textbf{Rob} & \textbf{3.5F} & \textbf{3F} & \textbf{3.1L} \\
\midrule
Fable 5    & \cmark & \xmark & \xmark & \xmark & \xmark & \xmark &
GPT-5.6-sol    & \cmark & \cmark & \cmark & \xmark & \xmark & \xmark &
Gemini 3.1 Pro      & \cmark & \cmark & \cmark & \cmark & \cmark & \cmark \\
Opus 4.8   & \cmark & \cmark & \cmark & \cmark & \cmark & \cmark &
GPT-5.6-terra  & \cmark & \cmark & \cmark & \xmark & \xmark & \xmark &
Gemini 3 Pro        & \cmark & \cmark & \cmark & \cmark & \cmark & \cmark \\
Sonnet 5   & \cmark & \cmark & \cmark & \cmark & \cmark & \cmark &
GPT-5.6-luna   & \cmark & \cmark & \cmark & \xmark & \xmark & \xmark &
Gemini Robotics 1.6  & \cmark & \cmark & \cmark & \cmark & \cmark & \cmark \\
Sonnet 4.6 & \cmark & \cmark & \cmark & \cmark & \cmark & \cmark &
GPT-5          & \cmark & \cmark & \cmark & \cmark & \xmark & \xmark &
Gemini 3.5 Flash    & \cmark & \cmark & \cmark & \cmark & \cmark & \cmark \\
Sonnet 4.5 & \cmark & \cmark & \cmark & \cmark & \cmark & \cmark &
GPT-5-mini     & \cmark & \cmark & \cmark & \xmark & \cmark & \xmark &
Gemini 3 Flash      & \cmark & \cmark & \cmark & \cmark & \cmark & \cmark \\
Haiku 4.5  & \cmark & \cmark & \cmark & \cmark & \cmark & \cmark &
o4-mini     & \cmark & \cmark & \cmark & \xmark & \xmark & \cmark &
Gemini 3.1 Flash Lite & \cmark & \cmark & \cmark & \cmark & \cmark & \cmark \\
\bottomrule
\end{tabular}%
}
\par\smallskip
\raggedright\scriptsize

\end{table}


\subsection{A Reasoning Extraction Attack}
\label{sec:extraction}

To extract reasoning, we exploit cross-session compatibility as described above. We port reasoning from a given target model, into the context window (see \Cref{fig:injection_schemes}) of a compatible but weaker model, which we coerce into transcribing the inserted reasoning token-by-token, using a simple ad-hoc jailbreak.

For each provider, we identify the weakest compatible model that enables extraction: Haiku 4.5 for Claude, as the weakest available model and one that supports assistant-turn prefilling; GPT-5.6 Luna for GPT, as the least capable model that interacts with reasoning traces from all earlier GPT models; and Gemini Robotics 1.6 for Gemini, as it can process traces from both the 2.5 and 3.x series (whereas we find that 3.1 Flash Lite does not interact with the reasoning from the 2.5 series). To measure the faithfulness of the extracted reasoning, we use the ratio of extracted reasoning tokens to API-reported thinking tokens. In this, we assume that API-based token counts are exact for billing reasons (at the time of writing), which allows us to use them as a form of ground-truth verification for total token count.

With this setup in place, we use the weaker models as ``fuzzy'' decoders for the encrypted reasoning blocks of stronger models with better safeguards. We provide a schematic illustration of this in \Cref{fig:teaser}. We detail the exact approach used
for each provider in \Cref{app:anthropic_extraction,app:openai_extraction,app:gemini_extraction}.

\paragraph{Faithfulness of the extracted reasoning.}
\label{sec:vuln_faithfulness}
In the absence of a ground-truth reasoning trace and the stochasticity of the generation process, we cannot guarantee that the extracted thoughts correspond exactly to
a model's private reasoning.  In \Cref{fig:teaser} we thus compare the API's reported thinking-token counts against the token counts of the extracted reasoning, re-encoded as input to the same model, on a set
of 120 Codeforces problems. For most of the inputs, the two counts track each other closely across all tested models, which is a good indicator for faithful extraction. To further verify correctness, we show in \cref{fig:distillation} that the extracted reasoning is indeed qualitatively more detailed than native summaries but also enables extraction of sensitive information otherwise not present in an input, such as API tokens and personally-identifiable information (\Cref{sec:av_secret_extraction}).

\paragraph{Why is this more scalable than directly jailbreaking a model?} A direct
extraction attack on a more capable model is
possible, but the attacker would have to bypass both model-level alignment (the
model's refusal to reveal its internal chain-of-thought) and system-level
defenses such as input filters and output substring-matching filters. The
availability of a less-capable model with compatible reasoning thus significantly lowers
the difficulty of successful extraction. Our experiments with Haiku 4.5, for instance, use a single fixed
extraction prompt across all attacks show in \Cref{fig:teaser}. By contrast, extraction from a relatively more capable GPT-5.6 Luna required different prompt templates for different reasoning blocks, best-of-$n$ sampling, and
ad-hoc workarounds for anti-distillation safeguards (e.g., splitting the extraction into chunks under 50 generated tokens).

\definecolor{usrbg}{RGB}{232,240,250}
\definecolor{usrfr}{RGB}{0,82,155}
\definecolor{thtbg}{RGB}{253,231,224}
\definecolor{thtfr}{RGB}{176,58,32}
\definecolor{txtfr}{RGB}{223,220,213}

\newcommand{\chipT}[1]{{\setlength{\fboxsep}{2.2pt}\setlength{\fboxrule}{.5pt}%
  \fcolorbox{thtfr}{thtbg}{\scriptsize\strut #1}}}
\newcommand{\chipX}[1]{{\setlength{\fboxsep}{2.2pt}\setlength{\fboxrule}{.5pt}%
  \fcolorbox{txtfr}{white}{\scriptsize\strut #1}}}

\begin{figure}[t]
\centering
\begin{tikzpicture}[
  turn/.style={rounded corners=2.5pt, draw=txtfr, line width=.6pt, fill=lstbg,
               inner xsep=4pt, inner ysep=4pt, font=\scriptsize},
  usr/.style={turn, fill=usrbg, draw=usrfr!55},
  arr/.style={-{Stealth[length=3.4pt,width=3pt]}, draw=black!42, line width=.6pt},
  ttl/.style={font=\scriptsize\bfseries, anchor=west},
  role/.style={font=\tiny, text=black!62, anchor=south},
  node distance=3.6mm,
]

\node[ttl] at (0,1.02) {Current Turn Injection};
\node[usr, anchor=west]  (c1) at (0,0) {user};
\node[turn, right=of c1] (c2) {\chipT{injected thought}};
\node[turn, right=of c2] (c3) {\chipX{text}};
\draw[arr] (c1) -- (c2);
\draw[arr] (c2) -- (c3);
\node[role] at (c2.north) {assistant};
\node[role] at (c3.north) {assistant continued};

\node[ttl] at ($(c3.east)+(1.5,1.02)$) {Past Turn Injection};
\node[usr, anchor=west]  (p1) at ($(c3.east)+(1.5,0)$) {user};
\node[turn, right=of p1] (p2) {\chipT{injected thought}\;\chipX{text}};
\node[usr,  right=of p2] (p3) {user};
\node[turn, right=of p3] (p4) {\chipX{text}};
\draw[arr] (p1) -- (p2);
\draw[arr] (p2) -- (p3);
\draw[arr] (p3) -- (p4);
\node[role] at (p2.north) {assistant};
\node[role] at (p4.north) {assistant};

\end{tikzpicture}
\looseness=-1 \caption{\textbf{Injection schemes exploiting in- and cross-session compatibility.} We find that
available form of thought injection depends on the provider and exact model. \textbf{Current-turn injection} places the thought in the current
assistant turn, so the model continues its visible answer straight from it; as of July 2026 it
is accepted by every GPT and Gemini model we tested and by the 4.5 generation of
Claude. In addition, Sonnet 4.5, Haiku 4.5, and all Gemini models accept
prefilling of the assistant turn's visible output. \textbf{Past-turn injection}
is available only against models that do not omit previous reasoning
blocks (e.g., Sonnet 5, Opus 4.8, Fable 5, and the GPT-5.6 series).}
\label{fig:injection_schemes}
\end{figure}


\subsection{Evaluation Setup}

All evaluations in this paper use the same basic extraction procedure: we obtain an encrypted reasoning block and replay it to a compatible decoder model using the provider-specific pipelines described in \Cref{app:anthropic_extraction,app:openai_extraction,app:gemini_extraction}. The evaluations differ only in how the encrypted reasoning blocks are collected.

\paragraph{Controlled extraction.}
For controlled experiments, we generate the reasoning blocks ourselves by querying target models through the Anthropic API, OpenAI API, and Google AI Studio. We use problems from AIME 2025 \citep{aime2025}, Codeforces (Open-R1 subset) \citep{codeforces}, and Humanity's Last Exam \citep{hle} to evaluate extraction fidelity and construct the qualitative examples presented throughout the paper. We also generate the source reasoning blocks ourselves for the jailbreaking and prompt-injection experiments.

\paragraph{Extraction from public traces.}
For the secret-extraction experiments, we instead collect raw agent and session traces published online by other users. We parse the encrypted reasoning blocks contained in these traces and decode them using the same extraction pipelines.

